\def\probno{PC}
\def\probtitle{SC ON 찾기}

\section{\probno{}. \probtitle{}}

\begin{frame}
    \sectiontitle{\probno{}}{\probtitle{}}
    \sectionmeta{
        \texttt{interactive, binary search}\\
        출제진 의도 -- \textbf{\color{acsilver}Medium}
    }
    \begin{itemize}
        \item 문제 아이디어: 조문성
        \item 문제 작업: 조문성
    \end{itemize}
\end{frame}

\begin{frame}{\probno{}. \probtitle{}}
    \begin{itemize}
        \item 질의한 집합 안에 보물이 하나라도 있으면 \texttt{1}, 없으면 \texttt{0}을 받습니다.
        \item 먼저 전체 위치 $[1,N]$에서 보물이 하나 있는 위치 $x$를 찾습니다.
        \item 구간의 왼쪽 절반을 질의합니다.
        \item 응답이 \texttt{1}이면 그 안에 보물이 있으므로 왼쪽으로 이동합니다.
        \item 응답이 \texttt{0}이면 오른쪽 절반에 보물이 있으므로 오른쪽으로 이동합니다.
    \end{itemize}
\end{frame}

\begin{frame}{\probno{}. \probtitle{}}
    \begin{itemize}
        \item 이렇게 하면 $O(\log N)$번의 질의로 보물 하나의 위치 $x$를 찾을 수 있습니다.
        \item 이제 $x$를 제외한 나머지 위치에서 같은 방식으로 이분탐색합니다.
        \item 질의할 때는 항상 $x$를 제외하고 물어봅니다.
        \item 그러면 두 번째 보물 $y$도 $O(\log N)$번의 질의로 찾을 수 있습니다.
        \item 총 질의 수는 $2\lceil \log_2 N\rceil \le 20$이므로 제한 $25$번 안에 충분합니다.
    \end{itemize}
\end{frame}

\begin{frame}{\probno{}. \probtitle{}}
    \begin{itemize}
        \item 정답은 \texttt{! x y} 형식으로 출력합니다.
        \item 모든 질의와 정답 출력 뒤에는 반드시 flush해야 합니다.
        \item 인터랙터는 비적응형이므로, 보물 위치는 질의 중 변하지 않습니다.
    \end{itemize}
\end{frame}